% !Rnw weave = knitr
% Edit 07-instrumental-variables.Rnw; the corresponding .tex is generated.
\documentclass[xcolor={dvipsnames},handout]{beamer}

\usepackage{assets/pres-template_MOW}
\usepackage{pdfpages}

%--------------------------------------------------------------------------
% Specific to this document ---------------------------------------
\usepackage{amsmath}
\usepackage{amssymb}

%--------------------------------------------------------------------------
% \setbeamercovered{transparent}

%%%%%%%%%%%%%%%%%%%%%%%%%%%%%%%%%%%
\setlength{\tabcolsep}{1.3pt}
\title{PLSC 48110/28110}
\subtitle{Lecture 7: Instrumental variables estimation}
\date{}
\author{Molly Offer-Westort}
\institute{Department of Political Science, \\University of Chicago}


\usepackage{assets/Sweave}
\renewcommand{\Sweavesize}{\scriptsize}
\renewcommand{\Routsize}{\scriptsize}
\hypersetup{pdftitle={PLSC 48110/28110 - Lecture 7: Instrumental variables estimation},pdfauthor={Molly Offer-Westort},pdfpagemode=UseNone}
\begin{document}



%-------------------------------------------------------------------------------%
\frame{\titlepage
\thispagestyle{empty}
}
%-------------------------------------------------------------------------------%
\begin{frame}{Why we like experiments}

\begin{itemize}
\item Identification comes for free: we can ensure that treatment assignment is independent of potential outcomes.
\item Methods for analysis and inference are (relatively) straightforward.
\item However, treatment assignment doesn't guarantee treatment uptake.  (examples: GOTV, de-worming; ethical considerations)
\item So we may have independence of treatment assignment and potential outcomes, but not \textit{actual} treatment received and potential outcomes.
\end{itemize}

\end{frame}


%%%%%NOTE%%%%%
\note{
\scriptsize \singlespacing

\begin{itemize}
\item xxxx
\end{itemize}

}

%-------------------------------------------------------------------------------%
\begin{frame}{Noncompliance}

\begin{itemize}
\item People who opt-out of treatment will often look different, and have different associated experimental responses, as compared to people who will opt-in when treatment is offered.
\item We can focus just on \textit{intention to treat} (ITT) effects: what is the effect on a group of having treatment offered to them  \textit{whether they take up treatment or not}.
\item Combination of effects among people who \textit{do} take up treatment when randomized to receive it, and (potentially null) effects among people who \textit{don't}.  (We'll come back to compliance types).
\item Often these are of policy relevance—for things you can't \textit{force} people to take, it's worth understanding what the average effect of making them available is.  And identification is still good.
\item However, this might not be a quantity that we're interested in, and when the assignment mechanism would look different in the ``real'' world, external validity may be poor.
\end{itemize}

\end{frame}


%%%%%NOTE%%%%%
\note{
\scriptsize \singlespacing

\begin{itemize}
\item xxxx
\end{itemize}

}

%-------------------------------------------------------------------------------%
\begin{frame}{Instrument}

\begin{itemize}
\item With noncompliance when we can't \textit{force} everyone to take treatment, we can consider effects on the subgroup for whom random assignment was effective at moving them to uptake treatment.
\item Our estimand becomes the "Local Average Treatment Effect" or "Complier Average Treatment Effect."
\end{itemize}

\end{frame}


%%%%%NOTE%%%%%
\note{
\scriptsize \singlespacing

\begin{itemize}
\item xxxx
\end{itemize}

}

%-------------------------------------------------------------------------------%
\begin{frame}{Observational data motivation}

\begin{itemize}
\item Motivation from natural experiments \citep{angrist2001instrumental}
\item Uses only a part of the variation in the treatment received
\item Relevance of estimand?  Do we think there is \textbf{heterogeneity} in response?
\item In observational setting, what is the alternative?
\end{itemize}

\end{frame}


%%%%%NOTE%%%%%
\note{
\scriptsize \singlespacing

\begin{itemize}
\item xxxx
\end{itemize}

}

%-------------------------------------------------------------------------------%
\begin{frame}{Instrument}

\begin{itemize}
\item Call the treatment \textit{assignment} an \textit{instrument} $Z_i\in \{0,1 \}$
\item Call the treatment \textit{received} $W_i$;
\begin{itemize}
\item Treatment \textit{received} is a function of treatment \textit{assigned}: $W_i(Z_i)\in \{0,1 \}$
\end{itemize}
\item We can denote  outcomes as a function of treatment assigned and received: $Y_i^{obs} = Y_i(Z_i, W_i(Z_i))$
% \item By random assignment, $Z_i \independent (Y_i(0), Y_i(1))$  (people \textit{assigned} to treatment look no different, on average, from people \textit{assigned} to control)
% \begin{itemize}
% \item (We can relax this to allow for independence conditional on covariates, for example if we're using stratified random assignment.)
% \end{itemize}
\end{itemize}

\end{frame}


%%%%%NOTE%%%%%
\note{
\scriptsize \singlespacing

\begin{itemize}
\item xxxx
\end{itemize}

}

%-------------------------------------------------------------------------------%
\begin{frame}{Compliance behavior}

$$
C_i =
\begin{cases}
\text{complier} & W_i(0) = 0,  W_i(1) = 1, \\
\text{defier} & W_i(0) = 1,  W_i(1) = 0, \\
\text{always-taker} & W_i(0) = 1,  W_i(1) = 1, \\
\text{never-taker} & W_i(0) = 0,  W_i(1) = 0
\end{cases}
$$

\end{frame}


%%%%%NOTE%%%%%
\note{
\scriptsize \singlespacing

\begin{itemize}
\item xxxx
\end{itemize}

}

%-------------------------------------------------------------------------------%
\begin{frame}{IV assumptions}

\begin{itemize}
\item No defiers (monotonicity): $W_i(1)\ge W_i(0)$  (is this plausible?)
\item Exclusion restriction: $Y_i(z,w) = Y_i(z',w), \text{ for all } z, z',w $  (is this plausible?)
\end{itemize}


$$
\Rightarrow Y_i(Z_i,W_i) = Y_i(W_i)
$$

\end{frame}


%%%%%NOTE%%%%%
\note{
\scriptsize \singlespacing

\begin{itemize}
\item xxxx
\end{itemize}

}

%-------------------------------------------------------------------------------%
\begin{frame}{ATE}

\begin{align*}
\tau^{ate} & = \E[Y_i(1) - Y_i(0)]
\end{align*}

But
\begin{align*}
\E[Y_i(1)] & \neq \E[Y_i^{obs}|Z_i = 1]\neq \E[Y_i^{obs}|W_i = 1]\\
\E[Y_i(0)] & \neq \E[Y_i^{obs}|Z_i = 0]\neq \E[Y_i^{obs}|W_i = 0]\\
\end{align*}


With noncompliance, ATE is not identified without further assumptions.

\end{frame}


%%%%%NOTE%%%%%
\note{
\scriptsize \singlespacing

\begin{itemize}
\item xxxx
\end{itemize}

}

%-------------------------------------------------------------------------------%
\begin{frame}{ITT}

\begin{align*}
\tau^{itt} & = \E[Y_i(1,W_i(1))] - \E[Y_i(0,W_i(0))] \\
            & = \E[Y_i(W_i)| Z_i = 1] - \E[Y_i(W_i)| Z_i = 0] \\
            & = \E[Y_i^{obs} | Z_i = 1] - \E[Y_i^{obs} | Z_i = 0]
\end{align*}


With noncompliance, ITT \textit{is} identified.


But it does not necessarily correspond to ATE.

\end{frame}


%%%%%NOTE%%%%%
\note{
\scriptsize \singlespacing

\begin{itemize}
\item xxxx
\end{itemize}

}

%-------------------------------------------------------------------------------%
\begin{frame}{LATE}

\begin{align*}
\tau^{late} & = \E[Y_i(1) - Y_i(0) | C_i = \text{complier}] \\
            & = \frac{\E[Y_i^{obs} | Z_i = 1] - \E[Y_i^{obs} | Z_i = 0]}{\E[W_i^{obs} | Z_i = 1] - \E[W_i^{obs} | Z_i = 0]}
\end{align*}



LATE also does not necessarily correspond to ATE!  Also different from ITT.

\end{frame}


%%%%%NOTE%%%%%
\note{
\scriptsize \singlespacing

\begin{itemize}
\item xxxx
\end{itemize}

}

%-------------------------------------------------------------------------------%
\begin{frame}[fragile]{Estimating methods}

Fabricated data: true treatment effect is 2

\begin{Schunk}
\begin{Sinput}
> library(fabricatr)
> library(estimatr)
> set.seed(60637)
> dat <- fabricate(
+   N = 1e4,
+   Z = rbinom(N, 1, prob = 0.4),
+   W  = Z * rbinom(N, 1, prob = 0.8),
+   Y = rpois(N, lambda = 4) + W * rpois(N, lambda = 2),
+   X = rnorm(N)
+ )
\end{Sinput}
\end{Schunk}



\end{frame}


%%%%%NOTE%%%%%
\note{
\scriptsize \singlespacing

\begin{itemize}
\item xxxx
\end{itemize}

}

%-------------------------------------------------------------------------------%
\begin{frame}[fragile]{Estimating methods}

\begin{Schunk}
\begin{Sinput}
> head(dat)
\end{Sinput}
\begin{Soutput}
     ID Z W  Y          X
1 00001 1 1 11 -0.7181921
2 00002 1 0  0  0.1173308
3 00003 1 1 11 -0.9540193
4 00004 1 1  6 -0.2415993
5 00005 0 0  5  1.0741809
6 00006 0 0  4  0.7435195
\end{Soutput}
\end{Schunk}



\end{frame}


%%%%%NOTE%%%%%
\note{
\scriptsize \singlespacing

\begin{itemize}
\item xxxx
\end{itemize}

}

%-------------------------------------------------------------------------------%
\begin{frame}[fragile]{Estimating methods: ITT}

Difference in means estimate
\begin{Schunk}
\begin{Sinput}
> mean(dat$Y[which(dat$Z == 1)]) -
+   mean(dat$Y[which(dat$Z == 0)])
\end{Sinput}
\begin{Soutput}
[1] 1.602577
\end{Soutput}
\end{Schunk}



\end{frame}


%%%%%NOTE%%%%%
\note{
\scriptsize \singlespacing

\begin{itemize}
\item xxxx
\end{itemize}

}

%-------------------------------------------------------------------------------%

\begin{frame}[fragile]{Estimating methods: LATE}

\begin{itemize}
\item Generically:
\begin{align*}
\hat \tau^{late} & = \widehat \Cov[Y_i, Z_i] / \widehat \Cov[W_i, Z_i]
\end{align*}
\item Wald estimator: With binary instrument, equivalently:
\begin{align*}
\hat \tau^{late}
            & = \frac{\hat \E[Y_i^{obs} | Z_i = 1] - \hat\E[Y_i^{obs} | Z_i = 0]}{\hat\E[W_i^{obs} | Z_i = 1] - \hat\E[W_i^{obs} | Z_i = 0]}
\end{align*}
\end{itemize}




\end{frame}


%%%%%NOTE%%%%%
\note{
\scriptsize \singlespacing

\begin{itemize}
\item xxxx
\end{itemize}

}

%-------------------------------------------------------------------------------%
\begin{frame}[fragile]{Estimating methods: LATE}

\small

\begin{Schunk}
\begin{Sinput}
> cov(dat$Y, dat$Z)/cov(dat$W, dat$Z)
\end{Sinput}
\begin{Soutput}
[1] 2.013196
\end{Soutput}
\end{Schunk}
\end{frame}

\begin{frame}[fragile]{Estimating methods: LATE (continued)}
\begin{Schunk}
\begin{Sinput}
> (numerator <- mean(dat$Y[which(dat$Z == 1)]) -
+    mean(dat$Y[which(dat$Z == 0)]))
\end{Sinput}
\begin{Soutput}
[1] 1.602577
\end{Soutput}
\begin{Sinput}
> (denominator <- mean(dat$W[which(dat$Z == 1)]) -
+     mean(dat$W[which(dat$Z == 0)]))
\end{Sinput}
\begin{Soutput}
[1] 0.7960361
\end{Soutput}
\begin{Sinput}
> numerator/denominator
\end{Sinput}
\begin{Soutput}
[1] 2.013196
\end{Soutput}
\end{Schunk}





\end{frame}


%%%%%NOTE%%%%%
\note{
\scriptsize \singlespacing

\begin{itemize}
\item xxxx
\end{itemize}

}

%-------------------------------------------------------------------------------%
\begin{frame}{Estimating methods: LATE}

\begin{itemize}
\item Two-stage least squares: regress treatment realized on instrument (and other instruments, covariates); regress outcome on predicted value of treatment (and control covariates)
\end{itemize}

\end{frame}


%%%%%NOTE%%%%%
\note{
\scriptsize \singlespacing

\begin{itemize}
\item xxxx
\end{itemize}

}

%-------------------------------------------------------------------------------%

\begin{frame}[fragile]{Estimating methods: LATE}


Doing each stage separately:
\footnotesize
\begin{Schunk}
\begin{Sinput}
> dat$W_predicted <- lm_robust(W ~Z, data = dat)$fitted.values
> lm_robust(Y ~ W_predicted, data = dat)
\end{Sinput}
\begin{Soutput}
            Estimate Std. Error   t value      Pr(>|t|) CI Lower
(Intercept) 4.008813 0.02590565 154.74667  0.000000e+00 3.958032
W_predicted 2.013196 0.05949577  33.83764 1.001255e-237 1.896573
            CI Upper   DF
(Intercept) 4.059593 9998
W_predicted 2.129820 9998
\end{Soutput}
\end{Schunk}

\end{frame}

\begin{frame}[fragile]{Estimating methods: LATE (continued)}
With \texttt{estimatr}:
\begin{Schunk}
\begin{Sinput}
> iv_robust(Y ~ W| Z, data = dat)
\end{Sinput}
\begin{Soutput}
            Estimate Std. Error  t value      Pr(>|t|) CI Lower
(Intercept) 4.008813 0.02590565 154.7467  0.000000e+00 3.958032
W           2.013196 0.05713043  35.2386 1.627323e-256 1.901209
            CI Upper   DF
(Intercept) 4.059593 9998
W           2.125184 9998
\end{Soutput}
\end{Schunk}


\end{frame}


%%%%%NOTE%%%%%
\note{
\scriptsize \singlespacing

\begin{itemize}
\item xxxx
\end{itemize}

}

%-------------------------------------------------------------------------------%
\begin{frame}[fragile]{Estimating methods: LATE}

With covariates:

\scriptsize
\begin{Schunk}
\begin{Sinput}
> iv_robust(Y ~ W + X | Z + X, data = dat)
\end{Sinput}
\begin{Soutput}
               Estimate Std. Error   t value      Pr(>|t|)
(Intercept)  4.00873400 0.02590631 154.73967  0.000000e+00
W            2.01258736 0.05712699  35.23006 2.134297e-256
X           -0.02730166 0.02131311  -1.28098  2.002305e-01
               CI Lower   CI Upper   DF
(Intercept)  3.95795242 4.05951558 9997
W            1.90060696 2.12456776 9997
X           -0.06907965 0.01447632 9997
\end{Soutput}
\end{Schunk}


\end{frame}


%%%%%NOTE%%%%%
\note{
\scriptsize \singlespacing

\begin{itemize}
\item xxxx
\end{itemize}

}

%-------------------------------------------------------------------------------%
\begin{frame}{Estimating methods: LATE}

\begin{itemize}
\item TSLS allows for multiple instruments, covariates
\item consistent but \textit{not} unbiased
\end{itemize}

\end{frame}


%%%%%NOTE%%%%%
\note{
\scriptsize \singlespacing

\begin{itemize}
\item xxxx
\end{itemize}

}

%-------------------------------------------------------------------------------%
\begin{frame}{Possible issues with instruments}

\begin{itemize}
\item Failure of exogeneity  (examples in the experimental setting?)
\item Weak instruments
\begin{itemize}
\item Diagnose strength using the first stage, $W \sim Z$ (plus any covariates): how strongly does the instrument predict treatment received? Use an appropriate first-stage $F$ statistic.
\item The reduced form, $Y \sim Z$, estimates the ITT effect; a small effect need not imply a weak instrument.
\end{itemize}
\end{itemize}

\end{frame}


%%%%%NOTE%%%%%
\note{
\scriptsize \singlespacing

\begin{itemize}
\item xxxx
\end{itemize}

}

%-------------------------------------------------------------------------------%
\begin{frame}{Alternatives}

\begin{itemize}
\item Assume unconfoundedness
$$
W_i \independent (Y_i(1), Y_i(0)) \rvert X_i
$$
\item Estimate average effect for compliers, generalize to population
\item Test equalities of compliers under treatment, always takers; compliers under control, never-takers; following \cite{athey2017econometrics}
\end{itemize}

\end{frame}


%%%%%NOTE%%%%%
\note{
\scriptsize \singlespacing

\begin{itemize}
\item xxxx
\end{itemize}

}

%-------------------------------------------------------------------------------%
\begin{frame}{Alternatives}

\begin{itemize}
\item Sometimes we can also \textit{bound} treatment effects.
\item Manski bounds: worst case bounds
\end{itemize}

\end{frame}


%%%%%NOTE%%%%%
\note{
\scriptsize \singlespacing

\begin{itemize}
\item xxxx
\end{itemize}

}

%-------------------------------------------------------------------------------%
\begin{frame}{Take-aways for experiments}

\begin{itemize}
\item What is the quantity of interest in your experimental setting?
\item When would you care about ATE? ITT? LATE?
\item Is non-compliance likely to be a problem?  (How do you define \textit{compliance}?)
\item Do you expect compliers and non-compliers to look different?
\item Do you expect compliers and non-compliers to have different (hypothetical) treatment effects?
\end{itemize}

\end{frame}


%%%%%NOTE%%%%%
\note{
\scriptsize \singlespacing

\begin{itemize}
\item xxxx
\end{itemize}

}

%-------------------------------------------------------------------------------%


\backupbegin
%-------------------------------------------------------------------------------%

\begin{frame}[allowframebreaks]{References}
    \bibliographystyle{apalike}
    \bibliography{assets/bib}
\end{frame}
%-------------------------------------------------------------------------------%
\backupend
\end{document}
